% Part: many-valued-logic
% Chapter: infinite-valued-logics
% Section: introduction

\documentclass[../../../include/open-logic-section]{subfiles}

\begin{document}

\olfileid{mvl}{inf}{int}

\olsection{引言}

一个矩阵的真值个数不必有限。对于无穷多真值构成的集合，一个显然的选择是 $0$ 与~$1$ 之间的有理数集，$V_\infty = [0,1] \cap \Rat$，
即：
\begin{align*}
    V_\infty & = \Setabs{\frac{n}{m}}{n,m \in \Nat \text{ 与 } n\le m}.
\intertext{在考察此无穷真值集时，常常也会考虑其子集：}
V_m & = \Setabs{\frac{n}{m-1}}{n \in \Nat \text{ 与 } n\le m}
\intertext{例如，$V_5$ 就是包含 $5$ 个等距真值的集合：}
V_5 & = \{0, \frac{1}{4}, \frac{1}{2}, \frac{3}{4}, 1\}.
\end{align*}
在基于这些真值集的逻辑中，通常只有 $1$ 是指定值，即 $V^+ = \{1\}$。
换句话说，我们令 $1$ 扮演（绝对）真的角色，$0$ 扮演绝对假的角色，但公式可以取~$V$ 中的任何中间值。

不过，我们也可以考虑所有 $0$ 与~$1$ 之间的\emph{实}数构成的集合 $V_{[0,1]} = [0,1]$，或者 $[0,1]$ 的其他无穷子集。
具有此类真值集的逻辑通常被称为\emph{模糊}逻辑。

\end{document}